\documentclass[12pt,a4paper]{article}
\usepackage[utf8]{inputenc}
\usepackage[T1]{fontenc}
\usepackage{amsmath,amssymb,amsfonts}
\usepackage{booktabs}
\usepackage{hyperref}
\usepackage[margin=2.5cm]{geometry}
\usepackage[numbers]{natbib}

\newcommand{\Lor}{\mathrm{Lor4D}}
\newcommand{\Icg}{I_{\mathrm{cg}}}

\title{Coarse-Graining Stability as a Cross-Scale Identity Test\\
  for Causal Set Dimension Selection}

\author{Gang Zhang}
\date{\today}

\begin{document}
\maketitle

\begin{abstract}
A robust dimension selection mechanism must produce an identity that survives
changes of resolution.
We test whether a 4D Lorentzian identity identified in finite causal-set
ensembles is preserved under coarse-graining.
Using a coarse-graining integrity index $\Icg$ that measures how well a poset
family's ranking survives node deletion, we establish a three-layer validation
protocol: (1)~identity retention under coarse-graining, (2)~ordering stability
across a penalty scan, and (3)~independent gain via block permutation.
In a frozen confirmatory window ($N=30$, keep ratio $0.6$, $\gamma=0.2$),
$\Icg$ predicts rank improvement with Spearman $\rho\approx 0.78$ ($p\leq
5\times 10^{-6}$, 200{,}000 permutations) across four independent replications.
Dose-response analysis confirms monotonic tertile ordering (9/9 conditions) with
all pooled $p<10^{-30}$.
A tie-aware Spearman fix (v12) resolves previously reported instability at
$N=52$, promoting it to a secondary confirmatory window ($\rho>0.55$, $p=5\times
10^{-6}$ across all variants and replications).
Perturbation quasi-intervention (5\%/10\%/20\% Hasse cover removal) confirms
dose-dependent CG disruption and near-perfect mechanistic coupling
($\rho\approx -0.98$ between $\Delta\Icg$ and $\Delta$penalty).
Within-stratum residualized tests recover layer-internal signal at all
perturbation doses after controlling for 11 confounders, upgrading the evidence
level from ``structural covariation'' to ``conditional quasi-causality.''
\end{abstract}

\noindent\textbf{Keywords:}
causal sets, coarse-graining, scale invariance, dimension selection,
structural stability

% ============================================================
\section{Introduction}
% ============================================================

Causal set theory models spacetime as a locally finite partially ordered set
with causal relations encoded by the order structure~\cite{Surya2019}.
A natural question is whether a selected 4D Lorentzian identity survives
changes of resolution:
if we coarse-grain a poset (remove a fraction of its elements), does the
$\Lor$ family retain its competitive advantage?

This question is not merely academic.
Any physical selection mechanism must be robust under renormalization-like
operations; an identity that dissolves at a different resolution is not a
genuine structural property but a scale-dependent artefact.

We address this through a coarse-graining (CG) stability analysis.
We define an integrity index $\Icg$ that measures how well a family's ranking
survives node deletion, and test whether $\Icg$ predicts rank improvement in a
three-layer validation protocol with pre-registered confirmatory windows.

% ============================================================
\section{Methods}
% ============================================================

\subsection{Coarse-graining operator}

Given an $N$-element poset $P$, we construct a coarse-grained version
$P_{\mathrm{cg}}$ by uniformly removing a fraction $(1-k)$ of elements
(keep ratio $k\in\{0.4,0.5,0.6,0.7,0.8\}$) and restricting the partial order
to the surviving elements.
The CG penalty measures three aspects of structural change:
drift (feature-space displacement), family switch (whether the nearest-family
label changes), and rank shift (change in competitive ranking).

\subsection{Integrity index}

$\Icg := -z(\text{penalty\_cg})$ within each block of fixed experimental
parameters.
Three variants are tested:
\textbf{full} (drift + switch + rank),
\textbf{no\_switch} (drift + rank, excluding the family-switch channel), and
\textbf{switch} (switch-only).

\subsection{Three-layer validation}

\textbf{Layer~1} (descriptive): Does $\Lor$ tend to retain family identity
after CG?

\textbf{Layer~2} (dynamic proxy): Is the ranking stable across a penalty scan?

\textbf{Layer~3} (confirmatory): Does $\Icg$ predict rank improvement
($\text{improve\_rank} = -\Delta\text{rank}$) under block permutation null?

% ============================================================
\section{Results}
% ============================================================

\subsection{Confirmatory window}

The frozen confirmatory window is $N=30$, keep ratio $0.6$, $\gamma=0.2$.
Four independent replications (seed offsets 172000/174000/175000/176000,
$\text{sis\_runs}=128$) yield consistent results (Table~\ref{tab:d_frozen}):

\begin{table}[h]
\centering
\caption{Confirmatory window results ($N=30$, $k=0.6$, $\gamma=0.2$).
$n_{\mathrm{perm}}=200{,}000$; $p\leq 5\times 10^{-6}$ denotes floor hit.}
\label{tab:d_frozen}
\begin{tabular}{lcccc}
\toprule
Variant & v10 $\rho$ & v12 $\rho$ & $p$ (v12) & Direction \\
\midrule
full       & 0.704 & 0.781 & $\leq 5\times 10^{-6}$ & $\uparrow$ \\
no\_switch & 0.501 & 0.716 & $\leq 5\times 10^{-6}$ & $\uparrow$ \\
switch     & 0.740 & 0.810 & $\leq 5\times 10^{-6}$ & $\uparrow$ \\
\bottomrule
\end{tabular}
\end{table}

The v12 tie-aware Spearman fix uniformly increases $\rho$ by resolving
artificial noise from tied integer ranks.

\subsection{N=52 promotion}

Previously reported instability at $N=52$ (rep5: no\_switch $\rho=0.096$,
$p=0.48$) was entirely a tie-handling artefact.
Under v12, all replications give $\rho>0.55$ and $p=5\times 10^{-6}$
(Table~\ref{tab:d_n52}):

\begin{table}[h]
\centering
\caption{$N=52$ under v12 tie-aware Spearman.}
\label{tab:d_n52}
\begin{tabular}{lccc}
\toprule
Rep & full $\rho$ & no\_switch $\rho$ & switch $\rho$ \\
\midrule
rep4 & 0.820 & 0.783 & 0.657 \\
rep5 & 0.615 & 0.553 & 0.564 \\
\bottomrule
\end{tabular}
\end{table}

\subsection{Dose-response}

Tertile binning of $\Icg$ (low/mid/high) shows monotonic increase in
improve\_rank across all 9 conditions (3 variants $\times$ 3 $N$ values).
All pooled $p<10^{-30}$.
Within-block OLS slopes are uniformly positive (36/36), with median Pearson
$r\approx 0.78$.

\subsection{Perturbation quasi-intervention}

Random removal of 5\%/10\%/20\% of Hasse cover relations, followed by
transitive closure reconstruction and full CG pipeline:

\begin{itemize}
\item Dose-strength confirmed: mean $|\Delta\Icg|$ monotonically increases
  with perturbation fraction (3/3 $N$ values).
\item Mechanistic coupling: $\Delta\Icg \leftrightarrow \Delta$penalty at
  $\rho\approx -0.98$ (sample-level, within-family).
\item Residualized within-stratum test (11 confounders): at 20\% perturbation,
  all 5 continuous $Y$ targets reach within-stratum significance
  ($p\leq 0.00002$), reversing the earlier ``layer-internal null'' finding.
\end{itemize}

\subsection{Cross-prediction check (C $\to$ D)}

Controlling for layer\_count and mean\_layer\_gap (Prediction~C indicators)
within the frozen window, $\Icg\to\text{improve\_rank}$ remains strongly
significant, confirming that CG stability is not merely a relabelling of
hierarchy depth.

% ============================================================
\section{Discussion}
% ============================================================

The D-line results establish that the $\Lor$ identity is not fragile under
resolution changes.
The evidence level is ``conditional quasi-causality'': after residualizing
within-stratum confounders, the perturbation-induced CG disruption predicts
changes in competitive fitness through mechanistically independent channels.

The dose-response monotonicity (9/9) and the near-perfect mechanistic coupling
($\rho\approx -0.98$) provide strong support for the interpretation that CG
stability reflects genuine structural quality rather than statistical artefact.

\subsection{Boundaries}

\begin{enumerate}
\item The confirmatory claim is a \emph{window claim} ($N=30$, $k=0.6$,
  $\gamma=0.2$), not a global theorem.
\item $\gamma=0.8$ is explicitly excluded (unstable under $\text{sis\_runs}=128$).
\item The drift-only variant shows sign flips and is not eligible for freezing.
\item The ``conditional quasi-causality'' label reflects that within-stratum
  signal requires residualization; raw stratified tests are null.
\end{enumerate}

% ============================================================
\section{Conclusion}
% ============================================================

We have demonstrated that the $\Lor$ identity is robust under coarse-graining
within a frozen confirmatory window, with dose-response monotonicity,
mechanistic coupling, and residualized within-stratum signal recovery.
This provides independent evidence that the finite-size Lorentzian identity is
structurally genuine, not a scale-dependent artefact.

\begin{thebibliography}{1}
\bibitem{Surya2019}
S.~Surya,
\emph{The causal set approach to quantum gravity},
\emph{Living Rev.\ Relativ.}\ \textbf{22}, 5 (2019),
arXiv:1903.11544.
\end{thebibliography}

\end{document}
